% !TEX program = xelatex
% hw: grouped numbering, label=, \inputproblem, problembreak, parts options,
% subparts, kind=, header=name-id, \akizumisetup{headleft=...}
\documentclass[kind=assignment, header=name-id, problembreak]{akizumi-hw}
\akizumisetup{course=Analysis, number=5, id=B12345678}
\author{Aaron Lin}
\newcommand{\expect}[2]{\edef\x{#1}\edef\y{#2}%
  \ifx\x\y\else\PackageError{test}{expected `#2', got `#1'}{}\fi}

\begin{document}
\maketitle
\makeatletter\expect{\@title}{Assignment\ 5}\makeatother

\begin{problem}[Warm-up]\label{prob:one}
Plain numbering still works.
\end{problem}
\expect{\theproblem}{1}

\problemset{B}
\inputproblem[Cauchy sequences]{hw/p1}
\inputproblem[title=Completeness]{hw/p2}
\expect{\theproblem}{B2}

\begin{problem}[label=B2$'$, title={Bonus, with a comma}]\label{prob:bonus}
A given label does not advance the count.
\end{problem}
\expect{\theproblem}{B2}

\problemset{C}
\begin{problem}[$f(x) = x^2$, a title with = and ,]
\begin{parts}[label=\arabic*.]
  \item first
  \item second
\end{parts}
\end{problem}
\expect{\theproblem}{C1}

\akizumisetup{headleft=Real Analysis, headright=Lin \textperiodcentered{} B12345678}
\begin{problem}
See \cref{prob:one}, \cref{prob:cauchy} and \cref{prob:bonus}.
\end{problem}

\end{document}
